\documentclass[11pt,a4paper,oneside]{report}
\usepackage[utf8]{inputenc}
\usepackage[T1]{fontenc}
\usepackage{lmodern}
\usepackage[polish]{babel}
\usepackage[margin=2.5cm]{geometry}
\usepackage{amsmath,amssymb,amsthm,mathtools}
\usepackage{array}
\usepackage{enumitem}
\usepackage{float}
\usepackage{tikz}
\usetikzlibrary{arrows.meta,calc,positioning,decorations.pathreplacing,decorations.markings,patterns}
\usepackage{pgfplots}
\pgfplotsset{compat=1.18}
\definecolor{manifoldblue}{RGB}{35,91,140}
\definecolor{manifoldred}{RGB}{170,65,55}
\definecolor{manifoldgreen}{RGB}{55,125,85}
\definecolor{manifoldgold}{RGB}{205,145,35}
\usepackage{microtype}
\usepackage{fancyhdr}
\usepackage[unicode]{hyperref}
\hypersetup{colorlinks=true, linkcolor=blue!45!black, citecolor=blue!45!black, urlcolor=blue!45!black}

\newcommand{\F}{\mathbb{F}}
\newcommand{\R}{\mathbb{R}}
\newcommand{\C}{\mathbb{C}}
\newcommand{\N}{\mathbb{N}}
\newcommand{\Z}{\mathbb{Z}}
\newcommand{\Q}{\mathbb{Q}}
\newcommand{\id}{\mathrm{id}}
\newcommand{\Lie}{\mathcal L}
\newcommand{\X}{\mathfrak X}
\newcommand{\supp}{\operatorname{supp}}
\newcommand{\pr}{\operatorname{pr}}
\newcommand{\D}{\mathrm D}
\newcommand{\dd}{\mathrm d}

\theoremstyle{plain}
\newtheorem{twierdzenie}{Twierdzenie}[section]
\newtheorem{lemat}[twierdzenie]{Lemat}
\newtheorem{stwierdzenie}[twierdzenie]{Stwierdzenie}
\newtheorem{wniosek}[twierdzenie]{Wniosek}
\newtheorem{fakt}[twierdzenie]{Fakt}
\theoremstyle{definition}
\newtheorem{definicja}[twierdzenie]{Definicja}
\newtheorem{przyklad}[twierdzenie]{Przykład}
\theoremstyle{remark}
\newtheorem{uwaga}[twierdzenie]{Uwaga}
\renewcommand{\proofname}{Dowód}

\pagestyle{fancy}
\fancyhf{}
\fancyhead[L]{\small Topologia algebraiczna}
\renewcommand{\chaptermark}[1]{\markboth{Rozdział \thechapter}{}}
\fancyhead[R]{\small\nouppercase{\leftmark}}
\fancyfoot[C]{\thepage}
\renewcommand{\headrulewidth}{0.4pt}

\setlength{\parindent}{0pt}
\setlength{\parskip}{4pt}
\setlength{\emergencystretch}{2em}

\begin{document}
\hypersetup{pageanchor=false}
\begin{titlepage}
\centering
\vspace*{4cm}
{\Huge\bfseries Topologia algebraiczna\par}
\vspace{0.6cm}
{\Large Homotopia, homologia, kohomologia i niezmienniki rozmaitości\par}
\vspace{1.3cm}
{\large Skrypt zebrany z rozdziałów książki o geometrii i topologii różniczkowej\par}
\vspace{1.0cm}
{\large Wersja XIX: po globalnym audycie matematycznym i redakcyjnym, 18 listopada 2026 r.\par}
\vfill
{\normalsize Wersja samodzielna; zadania w osobnym pliku.\par}
\end{titlepage}
\hypersetup{pageanchor=true}
\begingroup
\footnotesize\linespread{0.92}\selectfont\setlength{\parskip}{0pt}
\tableofcontents
\endgroup
\clearpage
\input{../rozdzialy/topologia_algebraiczna/00-wprowadzenie-i-zakres.tex}
\input{../rozdzialy/topologia_algebraiczna/01-jezyk-i-podstawowe-konstrukcje-topologiczne.tex}
\input{../rozdzialy/topologia_algebraiczna/02-podstawy-algebry-abstrakcyjnej-grupy-pierscienie-i-moduly.tex}
\input{../rozdzialy/topologia_algebraiczna/03-algebra-homologiczna-moduly-kompleksy-i-ko-homologie.tex}
\input{../rozdzialy/topologia_algebraiczna/04-homotopia-nakrycia-i-komorkowe-modele-przestrzeni.tex}
\input{../rozdzialy/topologia_algebraiczna/05-homologia-przestrzeni.tex}
\input{../rozdzialy/topologia_algebraiczna/06-wyzsze-grupy-homotopii-i-zwiazek-z-homologia.tex}
\input{../rozdzialy/topologia_algebraiczna/07-fibracje-i-dlugi-ciag-homotopii.tex}
\input{../rozdzialy/topologia_algebraiczna/08-kohomologia-singularna-i-iloczyn-kubkowy.tex}
\input{../rozdzialy/topologia_algebraiczna/09-formy-rozniczkowe-orientacja-i-kohomologia-de-rhama.tex}
\input{../rozdzialy/topologia_algebraiczna/10-homologia-i-kohomologia-c-echa-od-pokryc-do-granic-przestrzeni.tex}
\input{../rozdzialy/topologia_algebraiczna/11-przeciecia-dualnosc-i-klasy-charakterystyczne.tex}
\input{../rozdzialy/topologia_algebraiczna/12-ciagi-spektralne-od-filtracji-do-fibracji-serre-a.tex}
\input{../rozdzialy/topologia_algebraiczna/13-przestrzenie-klasyfikujace-i-teoria-przeszkod.tex}
\input{../rozdzialy/topologia_algebraiczna/14-wieze-postnikowa-iloczyny-whiteheada-i-operacje-steenroda.tex}
\input{../rozdzialy/topologia_algebraiczna/15-operacje-kohomologiczne-naturalnosc-iloczyny-i-relacje.tex}
\input{../rozdzialy/topologia_algebraiczna/16-zawieszenie-twierdzenie-freudenthala-i-homotopia-stabilna.tex}
\input{../rozdzialy/topologia_algebraiczna/17-konstrukcja-pontriagina-thoma-i-podstawy-teorii-kobordyzmow.tex}
\input{../rozdzialy/topologia_algebraiczna/18-homologia-morse-a-od-punktow-krytycznych-do-homologii-singular.tex}
\input{../rozdzialy/topologia_algebraiczna/19-stopien-indeks-punktu-stalego-i-twierdzenie-lefschetza.tex}
\input{../rozdzialy/topologia_algebraiczna/20-zwartosc-i-sklejanie-trajektorii-morse-a-podstawy-analityczne.tex}
\input{../rozdzialy/topologia_algebraiczna/21-indeks-pol-wektorowych-i-twierdzenie-poincarego-hopfa-z-brzegi.tex}
\input{../rozdzialy/topologia_algebraiczna/22-kohomologia-z-nosnikami-zwartymi-i-homologia-borel-moore-a.tex}
\input{../rozdzialy/topologia_algebraiczna/23-kohomologia-snopow-z-nosnikami-i-podstawy-dualnosci-verdiera.tex}
\input{../rozdzialy/topologia_algebraiczna/24-homologia-przeciec-snopy-perwersyjne-i-twierdzenie-o-rozkladzi.tex}
\input{../rozdzialy/topologia_algebraiczna/25-kategorie-pochodne-struktury-hodge-a-cykle-zanikajace-i-monodr.tex}
\input{../rozdzialy/topologia_algebraiczna/26-wariacje-struktur-hodge-a-koneksja-gaussa-manina-i-degeneracje.tex}
\input{../rozdzialy/topologia_algebraiczna/27-szesc-operacji-snopowych-dualnosc-verdiera-i-cykle-specjalizac.tex}
\input{../rozdzialy/topologia_algebraiczna/28-topologia-gladkich-zespolonych-odmian-projektowych-twierdzenia.tex}
\input{../rozdzialy/topologia_algebraiczna/29-dalsze-narzedzia-homotopii-granice-odwrotne-przeszkody-operacj.tex}
\input{../rozdzialy/topologia_algebraiczna/30-minimalne-przypomnienie-geometrii-rozniczkowej.tex}
\end{document}
